\appendixitem{PROJECTIVE LIMITS OF LIE GROUPS}

Lemma 1. Let \((\mathrm{G}_{\alpha}, f_{\alpha \beta})\) be a projective system of topological groups relative to a filtered index set I, and G its limit. Assume that the canonical maps \(f_{\alpha}: \mathrm{G} \to \mathrm{G}_{\alpha}\) are surjective.

\begin{enumerate}
\def\labelenumi{\alph{enumi})}
\item
  The subgroups \(\overline{\mathrm{D}(\mathrm{G}_{\alpha})}\) (resp. \(\mathrm{C}(\mathrm{G}_{\alpha})\) , resp. \(\mathrm{C}(\mathrm{G}_{\alpha})_0\) ) form a projective system of subsets of \(\mathrm{G}_{\alpha}\) .
\item
  We have \(\overline{\mathrm{D}(\mathrm{G})}=\varprojlim_{\alpha}\overline{\mathrm{D}(\mathrm{G}_{\alpha})}\) and \(\mathrm{C}(\mathrm{G})=\varprojlim_{\alpha}\mathrm{C}(\mathrm{G}_{\alpha})\) .
\item
  If \(\mathrm{G}_{\alpha}\) is compact for all \(\alpha \in \mathrm{I}\) , then \(\mathrm{C}(\mathrm{G})_0 = \varprojlim_{\alpha} \mathrm{C}(\mathrm{G}_{\alpha})_0\) .
\end{enumerate}

Let \(\alpha, \beta\) be two elements of I, with \(\alpha \leq \beta\) . Then \(f_{\alpha \beta}(\mathrm{D}(\mathrm{G}_{\beta})) \subset \mathrm{D}(\mathrm{G}_{\alpha})\) , and \(f_{\alpha \beta}(\mathrm{C}(\mathrm{G}_{\beta})) \subset \mathrm{C}(\mathrm{G}_{\alpha})\) since \(f_{\alpha \beta}\) is surjective; since \(f_{\alpha \beta}\) is continuous, it follows that \(f_{\alpha \beta}(\overline{\mathrm{D}(\mathrm{G}_{\beta})}) \subset \overline{\mathrm{D}(\mathrm{G}_{\alpha})}\) and \(f_{\alpha \beta}(\mathrm{C}(\mathrm{G}_{\beta})_0) \subset \mathrm{C}(\mathrm{G}_{\alpha})_0\) , hence \(a)\) . Since \(f_{\alpha}\) is surjective, \(f_{\alpha}(\mathrm{D}(\mathrm{G})) = \mathrm{D}(\mathrm{G}_{\alpha})\) (Algebra, Chap. I, §6, no. 2, Prop. 6), so \(\overline{\mathrm{D}(\mathrm{G})} = \varprojlim \overline{\mathrm{D}(\mathrm{G}_{\alpha})}\) (General Topology, Chap. I, §4, no. 4, Cor. of Prop. 9). The surjectivity of \(f_{\alpha}\) also implies the inclusion \(f_{\alpha}(\mathrm{C}(\mathrm{G})) \subset \mathrm{C}(\mathrm{G}_{\alpha})\) and hence \(\mathrm{C}(\mathrm{G}) \subset \varprojlim \mathrm{C}(\mathrm{G}_{\alpha})\) ; the opposite inclusion is immediate. Finally, assertion \(c)\) follows from \(b)\) and General Topology, Chap. III, §7, no. 2, Prop. 4).

Lemma 2. Let \((\mathrm{S}_a)_{a\in \mathrm{A}},(\mathrm{T}_b)_{b\in \mathrm{B}}\) be two finite families of almost simple, simply-connected Lie groups (Chap. III, § 9, no. 8, Def. 3), \(u:\prod_{a\in \mathrm{A}}\mathrm{S}_a\to \prod_{b\in \mathrm{B}}\mathrm{T}_b\) a surjective morphism. Then there exist an injective map \(l: \mathrm{B} \to \mathrm{A}\) and isomorphisms \(u_b: \mathrm{S}_{l(b)} \to \mathrm{T}_b\) ( \(b \in \mathrm{B}\) ) such that \(u((s_a)_{a \in \mathrm{A}}) = (u_b(s_{l(b)}))_{b \in \mathrm{B}}\) for every element \((s_a)_{a \in \mathrm{A}}\) of \(\prod_{a \in \mathrm{A}} \mathrm{S}_a\) .

Denote by \(\mathfrak{s}_a\) (resp. \(\mathfrak{t}_b\) ) the Lie algebra of \(S_a\) (resp. \(T_b\) ) for \(a \in A\) (resp. \(b \in B\) ), and consider the homomorphism \(L(u): \prod_{a \in A} \mathfrak{s}_a \to \prod_{b \in B} \mathfrak{t}_b\) . Its kernel is an ideal of the semi-simple Lie algebra \(\prod_{a \in A} \mathfrak{s}_a\) , and hence is of the form \(\prod_{a \in A''} \mathfrak{s}_a\) , with \(A'' \subset A\) (Chap. I, §6, no. 2, Cor. 1). Put \(A' = A - A''\) . By restriction, \(L(u)\) induces an isomorphism \(f: \prod_{a \in A'} \mathfrak{s}_a \to \prod_{b \in B} \mathfrak{t}_b\) . By loc. cit., for all \(a \in A'\) the ideal \(f(\mathfrak{s}_a)\) is equal to one of the \(\mathfrak{t}_b\) ; hence, there exists a bijection \(l: B \to A'\) such that \(f(\mathfrak{s}_{l(b)}) = \mathfrak{t}_b\) for \(b \in B\) , and \(f\) induces an isomorphism \(f_b: \mathfrak{s}_{l(b)} \to \mathfrak{t}_b\) . Since the groups \(S_a\) and \(T_b\) are simply-connected, there exist isomorphisms \(u_b: S_{l(b)} \to T_b\) such that \(L(u_b) = f_b\) for \(b \in B\) (Chap. III, §6, no. 3, Th. 3).

Denote by \(\tilde{u}:\prod_{a\in \mathrm{A}}\mathrm{S}_a\to \prod_{b\in \mathrm{B}}\mathrm{T}_b\) the morphism defined by \(\tilde{u} ((s_a)_{a\in \mathrm{A}}) = (u_b(s_{l(b)}))_{b\in \mathrm{B}}\) . By construction, \(\mathrm{L}(\tilde{u}) = f = \mathrm{L}(u)\) , so \(\tilde{u} = u\) , which proves the lemma.

Lemma 3. Under the hypotheses of Lemma 1, assume that the \(G_{\alpha}\) are simply-connected compact Lie groups. Then, the topological group G is isomorphic to the product of a family of almost simple, simply-connected compact Lie groups.

For all \(\alpha \in \mathrm{I}\) , the group \(\mathrm{G}_{\alpha}\) is the direct product of a finite family of almost simple, simply-connected subgroups \((\mathrm{S}_{\alpha}^{\lambda})_{\lambda \in \mathrm{L}_{\alpha}}\) (Chap. III, §9, no. 8, Prop. 28). Let \(\beta \in \mathrm{I}\) , \(\beta \geq \alpha\) . By Lemma 2, there exists a map \(l_{\beta \alpha}: \mathrm{L}_{\alpha} \to \mathrm{L}_{\beta}\) such that \(f_{\alpha \beta}(\mathrm{S}_{\beta}^{l_{\beta \alpha}(\lambda)}) = \mathrm{S}_{\alpha}^{\lambda}\) for \(\lambda \in \mathrm{L}_{\alpha}\) . We have \(l_{\gamma \beta} \circ l_{\beta \alpha} = l_{\gamma \alpha}\) for \(\alpha \leq \beta \leq \gamma\) , so \((\mathrm{L}_{\alpha}, l_{\beta \alpha})\) is an inductive system of sets relative to I. Let L be its limit; the maps \(l_{\beta \alpha}\) being injective, \(\mathrm{L}_{\alpha}\) can be identified with a subset of L, so that \(\mathrm{L} = \bigcup_{\alpha \in \mathrm{I}} \mathrm{L}_{\alpha}\) .

Let \(\lambda \in L\) . Put \(S_{\alpha}^{\lambda} = \{1\}\) when \(\lambda \notin L_{\alpha}\) , and denote by \(\varphi_{\alpha \beta}^{\lambda}: S_{\beta}^{\lambda} \to S_{\alpha}^{\lambda}\) the morphism induced by \(f_{\alpha \beta}\) ; this gives a projective system of topological groups \((S_{\alpha}^{\lambda}, \varphi_{\alpha \beta}^{\lambda})\) , whose limit is isomorphic to \(S_{\lambda}\) . The canonical homomorphism of topological groups

\[
\lim_{\substack{\longleftarrow \\ \alpha \in I}}\left(\prod_{\lambda \in L}S_{\alpha}^{\lambda}\right)\to \prod_{\lambda \in L}\left(\lim_{\substack{\longleftarrow \\ \alpha \in I}}S_{\alpha}^{\lambda}\right)
\]

is bijective (Theory of Sets, Chap. III, §7, no. 3, Cor. 2); it is thus an isomorphism since the groups in question are compact. But the first of these groups can be identified with G and the second with the product of the \(S_{\lambda}\) , hence the lemma.

